\documentclass[11pt]{article}

\usepackage[a4paper, margin=1in]{geometry}
\usepackage{amsmath, amssymb, amsthm}
\usepackage{mathtools}
\usepackage{hyperref}
\usepackage{enumitem}
\usepackage{parskip}

\title{Why a Time-Dependent Neural-Network Potential Correction\\
Cannot Recover $h_{\rm true}$ Under an Initial-PDF Mismatch}
\author{}
\date{}

\newtheorem*{claim}{Claim}

\begin{document}
\maketitle

\section*{Setup}

We back-integrate observed stellar coordinates $(z_i, w_i)$ at observation
time $t_{\rm obs}$ to initial coordinates $(z_i(0), w_i(0))$ using a model
gravitational potential
\begin{equation}
  \Phi_{\rm model}(z, t) \;=\; \Phi_{\rm static}(z;\,h) \;+\; V_{\rm nn}(z, t),
\end{equation}
where $\Phi_{\rm static}$ is the fixed baryon $+$ dark-matter-halo contribution
parametrised by the dark-matter scale height $h$, and $V_{\rm nn}(z, t)$ is an
arbitrary neural-network potential correction.  The likelihood assumes a
Gaussian initial distribution,
\begin{equation}
  f_0(z, w) \;=\; \mathcal{N}(0, \sigma_{z,\rm fit})\;\times\;\mathcal{N}(0, \sigma_{w,\rm fit}),
\end{equation}
and the maximum-likelihood estimate minimises
\begin{equation}
  \mathrm{NLL}(h, V_{\rm nn}) \;=\; \sum_i \left[\,\tfrac{1}{2}\,\frac{z_i(0)^2}{\sigma_{z,\rm fit}^2} + \tfrac{1}{2}\,\frac{w_i(0)^2}{\sigma_{w,\rm fit}^2}\,\right] + \mathrm{const}.
\label{eq:nll}
\end{equation}
over $(h, V_{\rm nn})$ jointly.  The data was generated with
$(\sigma_{z,\rm true}, \sigma_{w,\rm true})$.  The ``PDF mismatch'' we study is
the case in which $\sigma_{w,\rm fit} < \sigma_{w,\rm true}$ (the argument is
identical for $\sigma_{z,\rm fit} < \sigma_{z,\rm true}$).

\section*{Claim}

\begin{claim}
No $V_{\rm nn}(z, t)$ --- regardless of architecture, time-dependence, network
capacity, or physical constraints (positivity, smoothness, finite mass, etc.)
--- can move the joint MLE to $h^* = h_{\rm true}$.
\end{claim}

The argument below proceeds in five short steps.  The conclusion is
\emph{structural}: the bias lives in the loss function itself, not in the model
space, and no enrichment of the model space can compensate for that.

\section*{Step 1: At the truth, the loss is large, not small}

Plug in the correct model: $h = h_{\rm true}$, $V_{\rm nn} \equiv 0$.  The
dynamics are then exactly those used to generate the data, and the
back-integration is exact:
\begin{equation}
  w_i(0) \sim \mathcal{N}(0,\, \sigma_{w,\rm true}),
  \qquad
  z_i(0) \sim \mathcal{N}(0,\, \sigma_{z,\rm true}).
\end{equation}
The expected per-star $w$-term in the NLL is
\begin{equation}
  \mathbb{E}\!\left[\,\tfrac{1}{2}\,\frac{w(0)^2}{\sigma_{w,\rm fit}^2}\,\right]
  \;=\;
  \tfrac{1}{2}\,\frac{\sigma_{w,\rm true}^2}{\sigma_{w,\rm fit}^2}.
\end{equation}
When $\sigma_{w,\rm fit} < \sigma_{w,\rm true}$, this is much larger than the
$\tfrac{1}{2}$ that would obtain if the PDF were correct.  The truth is
\emph{not} a stationary point of the loss; the gradient
$\partial\mathrm{NLL}/\partial h$ at $(h_{\rm true}, 0)$ is non-zero, and the
optimiser will move \emph{away} from the truth in order to reduce the loss.

\section*{Step 2: The loss minimum lies wherever the back-integrated\\ $w$-distribution is narrow}

The optimiser wants the back-integrated initial-velocity variance to satisfy
\begin{equation}
  \langle w(0)^2 \rangle \;\approx\; \sigma_{w,\rm fit}^2
\end{equation}
so that the squared term in~\eqref{eq:nll} is of order one.  Two
\emph{physically independent} mechanisms in the model can produce such a narrow
$w(0)$ distribution:
\begin{description}[leftmargin=2em, style=nextline]
  \item[(a) Drift in $h$.]
  Lowering $h$ increases the magnitude of the anharmonic $z^4/h^2$ correction
  in $\Phi_{\rm DM}(z)$.  High-amplitude orbits become ``stiffer'' --- their
  back-integration compresses them more --- so the distribution of $w(0)$ is
  narrower.

  \item[(b) A positive-curvature $V_{\rm nn}$.]
  Adding a midplane mass concentration produces a quadratic correction
  $V_{\rm nn}(z) \approx \tfrac{1}{2}\,\Delta\omega^2\, z^2$ near the midplane.
  This raises the effective harmonic frequency at $z = 0$, which in turn
  modifies the back-integration so that $w(0)$ is narrower.  A flexible NN
  finds this very easily.
\end{description}

The MLE landscape is therefore \emph{not} point-like at any single minimum; it
has a continuous ridge of nearly-degenerate solutions parametrised by
trade-offs between (a) and (b).  The truth $(h_{\rm true}, 0)$ lies \emph{off}
this ridge.  Adam picks a point on the ridge based on initialisation and
learning rates, but it will not, in general, land at the truth.

\section*{Step 3: Time-dependence cannot help --- it can only enlarge\\
the degeneracy}

A time-dependent correction $V_{\rm nn}(z, t)$ is a strict superset of
$V_{\rm nn}(z)$:
\begin{itemize}[leftmargin=2em]
  \item Every static $V_{\rm nn}$ that absorbs the PDF bias is still available
        as a $t$-constant member of the larger class.
  \item Additional time-dependent modes --- for example, periodic stiffening of
        the well at orbital frequencies --- can produce \emph{additional}
        mechanisms that narrow the back-integrated $w(0)$ distribution.
\end{itemize}

The dimension of the degenerate ridge grows.  Time-dependence \emph{cannot
remove} any direction; it can only add more low-NLL configurations to choose
from.  A time-dependent NN therefore cannot do better than the static NN --- it
can only do equally badly or worse.

\section*{Step 4: Physical constraints (e.g.\ $\rho_{\rm nn} \ge 0$ via Poisson)\\
do not remove the relevant ridge}

A natural physical restriction is to require that the NN's potential
correspond to a non-negative mass distribution:
\begin{equation}
  \frac{\partial^2 V_{\rm nn}}{\partial z^2} \;=\; 4\pi G\,\rho_{\rm nn}(z) \;\ge\; 0.
\end{equation}
This forces $V_{\rm nn}$ to be everywhere convex.  It successfully blocks one
direction of bias-absorption --- concave corrections like $-z^4/h^2$ that would
mimic $h$-decreases by reshaping the anharmonic structure of the potential.

But it does \emph{not} block mechanism (b).  A constant positive density near
midplane gives a quadratic $V_{\rm nn}(z) \approx \mathrm{const}\cdot z^2$,
which is convex --- fully allowed under $\rho_{\rm nn} \ge 0$.  This is exactly
the midplane-mass-concentration that narrows $w(0)$.  The bias-absorbing
direction therefore \emph{survives} the positivity constraint.  The ridge
persists.

Any constraint we can write down on the dynamics-side model can only carve
along the loss surface.  It cannot move the surface itself to a place where
$(h_{\rm true}, 0)$ is the optimum, because that is not what the constraint
does.

\section*{Step 5: The bias is in the loss function, not the model}

The structural obstruction is this.  The loss
\begin{equation}
  \mathrm{NLL}(h, V_{\rm nn}) \;\sim\; \sum_i \frac{w_i(0;\,h, V_{\rm nn})^2}{\sigma_{w,\rm fit}^2} + \cdots
\end{equation}
has $\sigma_{w,\rm fit}$ \textbf{frozen at the wrong value}.  The minimum of
any such loss over any model class --- potential corrections, time-dependent
corrections, physics-informed corrections, neural-network corrections, all of
them --- sits at the model point that produces
$\langle w(0)^2\rangle \approx \sigma_{w,\rm fit}^2$.  By assumption that is
not the truth, because the truth produces
$\langle w(0)^2\rangle \approx \sigma_{w,\rm true}^2 > \sigma_{w,\rm fit}^2$.

\begin{quote}
\emph{No flexibility on the model side can change where the loss function's
minimum sits, because the loss function is defined by the wrong
$\sigma_{w,\rm fit}$.}
\end{quote}

A NN potential correction is a tool for redistributing flexibility
\emph{within} the dynamics.  The bias does not live in the dynamics; it lives
in the loss function, via the assumed initial distribution.  The tool is
pointed at the wrong target.

\section*{Corollary: What \emph{can} recover $h_{\rm true}$}

Only two strategies can succeed.

\subsection*{(1) Change the loss function}

Make $\sigma_{w,\rm fit}$ (and $\sigma_{z,\rm fit}$) \textbf{free parameters}
of the fit.  The joint MLE then minimises over
$(h, \sigma_{z,\rm fit}, \sigma_{w,\rm fit}, \ldots)$, and the joint minimum
is at $(h_{\rm true}, \sigma_{z,\rm true}, \sigma_{w,\rm true}, \ldots)$.  This
is the structurally correct fix: the bias was that the loss had a wrong frozen
constant; allowing that constant to float removes the bias at its source.

\subsection*{(2) Add information beyond the per-star NLL}

Cross-validation against held-out stars, multiple datasets at different
$t_{\rm obs}$, an external prior on $h$, observed spatial structure of
$\sigma_w(z)$, etc.  Each of these adds a constraint that the wrong-$\sigma_w$
bias cannot satisfy simultaneously, because the bias produces different
signatures at different datasets or scales.

\subsection*{Where a NN can usefully live}

A NN solution to the PDF-mismatch problem is to put NN flexibility where the
bias actually lives --- in the initial PDF itself.  Replace the rigid Gaussian
$f_0(z, w)$ with a NN-parametrised density (e.g.\ a normalising flow over
$(z, w)$, or its parametric special case: a Gaussian whose parameters are
learnable, or a Gaussian mixture).  The likelihood becomes
\begin{equation}
  \mathrm{NLL} \;=\; -\sum_i \log f_{0,\rm NN}\bigl(z_i(0), w_i(0)\bigr).
\end{equation}
This NN can absorb the PDF mismatch but \emph{cannot shift orbits}, so it has
no leverage on $h$ other than the structural information in the
back-integration.  The joint MLE then recovers $h \approx h_{\rm true}$
together with the correct $f_0$.

\section*{One-sentence summary}

A NN potential correction adds knobs to the \emph{dynamics}, but the bias from
a wrong initial PDF lives in the \emph{loss function} --- so the dynamics-side
knobs can never reach the truth, while a NN that parametrises the initial PDF
itself can.

\end{document}
